% Options for packages loaded elsewhere
% Options for packages loaded elsewhere
\PassOptionsToPackage{unicode}{hyperref}
\PassOptionsToPackage{hyphens}{url}
\PassOptionsToPackage{dvipsnames,svgnames,x11names}{xcolor}
%
\documentclass[
  letterpaper,
  DIV=11,
  numbers=noendperiod]{scrartcl}
\usepackage{xcolor}
\usepackage{amsmath,amssymb}
\setcounter{secnumdepth}{-\maxdimen} % remove section numbering
\usepackage{iftex}
\ifPDFTeX
  \usepackage[T1]{fontenc}
  \usepackage[utf8]{inputenc}
  \usepackage{textcomp} % provide euro and other symbols
\else % if luatex or xetex
  \usepackage{unicode-math} % this also loads fontspec
  \defaultfontfeatures{Scale=MatchLowercase}
  \defaultfontfeatures[\rmfamily]{Ligatures=TeX,Scale=1}
\fi
\usepackage{lmodern}
\ifPDFTeX\else
  % xetex/luatex font selection
\fi
% Use upquote if available, for straight quotes in verbatim environments
\IfFileExists{upquote.sty}{\usepackage{upquote}}{}
\IfFileExists{microtype.sty}{% use microtype if available
  \usepackage[]{microtype}
  \UseMicrotypeSet[protrusion]{basicmath} % disable protrusion for tt fonts
}{}
\makeatletter
\@ifundefined{KOMAClassName}{% if non-KOMA class
  \IfFileExists{parskip.sty}{%
    \usepackage{parskip}
  }{% else
    \setlength{\parindent}{0pt}
    \setlength{\parskip}{6pt plus 2pt minus 1pt}}
}{% if KOMA class
  \KOMAoptions{parskip=half}}
\makeatother
% Make \paragraph and \subparagraph free-standing
\makeatletter
\ifx\paragraph\undefined\else
  \let\oldparagraph\paragraph
  \renewcommand{\paragraph}{
    \@ifstar
      \xxxParagraphStar
      \xxxParagraphNoStar
  }
  \newcommand{\xxxParagraphStar}[1]{\oldparagraph*{#1}\mbox{}}
  \newcommand{\xxxParagraphNoStar}[1]{\oldparagraph{#1}\mbox{}}
\fi
\ifx\subparagraph\undefined\else
  \let\oldsubparagraph\subparagraph
  \renewcommand{\subparagraph}{
    \@ifstar
      \xxxSubParagraphStar
      \xxxSubParagraphNoStar
  }
  \newcommand{\xxxSubParagraphStar}[1]{\oldsubparagraph*{#1}\mbox{}}
  \newcommand{\xxxSubParagraphNoStar}[1]{\oldsubparagraph{#1}\mbox{}}
\fi
\makeatother


\usepackage{longtable,booktabs,array}
\newcounter{none} % for unnumbered tables
\usepackage{calc} % for calculating minipage widths
% Correct order of tables after \paragraph or \subparagraph
\usepackage{etoolbox}
\makeatletter
\patchcmd\longtable{\par}{\if@noskipsec\mbox{}\fi\par}{}{}
\makeatother
% Allow footnotes in longtable head/foot
\IfFileExists{footnotehyper.sty}{\usepackage{footnotehyper}}{\usepackage{footnote}}
\makesavenoteenv{longtable}
\usepackage{graphicx}
\makeatletter
\newsavebox\pandoc@box
\newcommand*\pandocbounded[1]{% scales image to fit in text height/width
  \sbox\pandoc@box{#1}%
  \Gscale@div\@tempa{\textheight}{\dimexpr\ht\pandoc@box+\dp\pandoc@box\relax}%
  \Gscale@div\@tempb{\linewidth}{\wd\pandoc@box}%
  \ifdim\@tempb\p@<\@tempa\p@\let\@tempa\@tempb\fi% select the smaller of both
  \ifdim\@tempa\p@<\p@\scalebox{\@tempa}{\usebox\pandoc@box}%
  \else\usebox{\pandoc@box}%
  \fi%
}
% Set default figure placement to htbp
\def\fps@figure{htbp}
\makeatother





\setlength{\emergencystretch}{3em} % prevent overfull lines

\providecommand{\tightlist}{%
  \setlength{\itemsep}{0pt}\setlength{\parskip}{0pt}}



 


\KOMAoption{captions}{tableheading}
\makeatletter
\@ifpackageloaded{caption}{}{\usepackage{caption}}
\AtBeginDocument{%
\ifdefined\contentsname
  \renewcommand*\contentsname{Table of contents}
\else
  \newcommand\contentsname{Table of contents}
\fi
\ifdefined\listfigurename
  \renewcommand*\listfigurename{List of Figures}
\else
  \newcommand\listfigurename{List of Figures}
\fi
\ifdefined\listtablename
  \renewcommand*\listtablename{List of Tables}
\else
  \newcommand\listtablename{List of Tables}
\fi
\ifdefined\figurename
  \renewcommand*\figurename{Figure}
\else
  \newcommand\figurename{Figure}
\fi
\ifdefined\tablename
  \renewcommand*\tablename{Table}
\else
  \newcommand\tablename{Table}
\fi
}
\@ifpackageloaded{float}{}{\usepackage{float}}
\floatstyle{ruled}
\@ifundefined{c@chapter}{\newfloat{codelisting}{h}{lop}}{\newfloat{codelisting}{h}{lop}[chapter]}
\floatname{codelisting}{Listing}
\newcommand*\listoflistings{\listof{codelisting}{List of Listings}}
\makeatother
\makeatletter
\makeatother
\makeatletter
\@ifpackageloaded{caption}{}{\usepackage{caption}}
\@ifpackageloaded{subcaption}{}{\usepackage{subcaption}}
\makeatother
\usepackage{bookmark}
\IfFileExists{xurl.sty}{\usepackage{xurl}}{} % add URL line breaks if available
\urlstyle{same}
% fallback for those not using the hyperref driver hyperxmp:
\makeatletter
\@ifundefined{xmpquote}{\newcommand{\xmpquote}[1]{#1}}{}
\makeatother
\hypersetup{
  pdftitle={Research Statement},
  pdfauthor={Montaque Reynolds, Phd},
  colorlinks=true,
  linkcolor={blue},
  filecolor={Maroon},
  citecolor={Blue},
  urlcolor={Blue},
  pdfcreator={LaTeX via pandoc}}


\title{Research Statement}
\author{Montaque Reynolds, Phd}
\date{}
\begin{document}
\maketitle


I work in areas of moral and experimental epistemology, primarily with
special concern for the relationship between emotion and knowledge. I
have published work on the expression of sentiment in popular media,
drawing attention to constructions of social identity that entail
relational agency through love and regret. I have work purporting to
uncover ascriptions of relationships that exist between members of
represented characters. By determining the valence, emotions, and other
affectual attributions from music lyrics, I have suggested that a more
comprehensive project in this area will go further towards determining
the portrayed broader cultural attitudes about inter-relational
dispositions of a given audience.

I began my PhD studies in 2018. In 2019 and 2020, I was the assistant to
the editor of the department's journal, Res Philosophica. Since its
founding in 1925 (as The Schoolman), Res Philosophica has been a
particularly important publication in academic philosophy. From 2020 to
2022 I was also the assistant to the editor of Schutzian Research, a
central publication dedicated to the thought of Alfred Schutz. In both
positions, I streamlined the editorial process, enhancing publication
efficiency. In part, I adopted an algorithm which extracted bibliography
information and encoded it according to latex standards. I also
conducted thorough manuscript reviews, ensuring high-quality content.
For each issue, I coordinated with authors, reviewers, and editorial
board to facilitate a smooth publication process. This enabled me to
foster strong relationships and improve communication and productivity.
I worked to ensure the authors adhered to submission deadlines, ensuring
timely publications.

From 2022 to 2023 I was the graduate assistant for the Walter J. Ong
Digital Humanities Center. While there, I conducted digital humanities
research, analyzing historical texts using computational methods and
techniques from natural language processing. I helped students and
faculty develop digital humanities projects. In at least one example
project, I helped students design projects analyzing and recovering
expressed sentiment from a group of texts considered formative to modern
Hindu culture. We built a database to store student annotations of the
materials and a website to host the findings and developed data
visualization tools to present findings effectively. I also pioneered
novel approaches in digital text analysis, integrating machine learning
algorithms to uncover hidden patterns in literary works. I often
partnered with faculty to design and implement digital projects,
enhancing interdisciplinary research capabilities across humanities
departments. This required that I meticulously curate and maintain
digital archives, ensuring data integrity and accessibility for future
research endeavors. I also applied computational methods to analyze
historical and contemporary texts.

I am currently a member of the Saint Louis Digital Humanities Network.
This is a group of educators in the Saint Louis area that are committed
to exploring ways and implications of expanding use of technology for
pedagogical purposes. As our digital technology develops, more and more
of our social lives are being conducted through digital mediums. Given
that my philosophical interests are motivated by questions about our
social environments, these new modes of social interaction are of great
interest. I am currently working on developing a philosophy course that
I intend to conduct in virtual environments, namely in Massive
Multiplayer Online communities (MMOs). I hope to collectively explore
implications for conducting our lives in digitally represented social
environments by using philosophical reflection to better understand how
online communities can be used as a social learning environment.

I have published work on adaptive preferences with members of CUNY's
graduate center.
https://www.taylorfrancis.com/chapters/edit/10.4324/9781003301424-19/adaptive-preferences-urna-chakrabarty-romy-feiertag-anne-marie-mccallion-brian-mcniff-jesse-prinz-montaque-reynolds-sukhvinder-shahi-maya-von-ziegesar-angella-yamamoto.
Additionally, I have 3 projects in draft states: (2025). ``Social Agency
and other Epistemic Scaffolding''; ``Deferrence and The Moral Properties
of Imaginative Experience in Narratives''; and ``Well-Being and Moral
Worth'' which explores semantic expressions of well-being in represented
human relationships. I also currently have a book titled Play and
Meaning, an epistemology of emotional testimony in contract with Vernon
Publishers. In it, I explore the view that stories are a form of play
that use our emotions to teach us what it means to live well. My current
research explores concepts related to imagination and well-being.

During the spring semester 2026, I began using Dave Chalmers' book
Reality+ in my intro classes. This book is primarily an exploration of
the philosophical method as applied not only to a single question, the
simulation hypothesis, but also a way that different areas of philosophy
may attempt to answer this question and why these attempts matter for
wider academic knowledge. Since then, I have become more interested in
the application of this work inside simulated environments. Using
various simulated environments, for instance, the Open world Space Sim
game Elite Dangerous, I round that my students were able to recognize
the gravity of philosophical reflection. Currently, my work is focusing
on virtual and simulated environments, for instance, the Classical
Ballet, exploring the possibility of works of art inside the space sim
Star Citizen, or using embodied LLM agents in Minecraft to better focus
on the importance of epistemic justice and social relationships for
higher order learning.

I am excited about the opportunity to contribute to a learning
environment that engages the student population in meaningful
philosophical discourse. My interdisciplinary experience, and
philosophical rigor combined with innovative pedagogical methods will be
a terrific addition to your department. I would welcome the opportunity
to discuss my qualifications and the needs of your department further.
Thank you for your time, and for your thoughtful consideration.




\end{document}
